\section{What John 6 Actually Does}

The chapter does not begin with a sacramental formula. It begins with hungry
people and a sign. Jesus takes five barley loaves, gives thanks, distributes,
and leaves a superabundance (Jn 6:1--15). The crowd reads the sign politically
and seeks to make him king. After the crossing of the sea, Jesus exposes a
narrower appetite: they seek him because they ate and were filled, not because
they understood the sign (6:26).

The discourse then advances by repeated corrections. The crowd invokes the
manna. Jesus distinguishes the gift, the giver, and the life bestowed: Moses
did not give the true bread; the Father gives it; the bread of God comes down
from heaven and gives life to the world (6:31--33). When they ask for this
bread, he answers with a person rather than an object: ``I am the bread of
life'' (6:35). Coming and believing answer hunger and thirst. The first movement
therefore cannot be reduced to the mechanics of eating. Faith in the one sent
by the Father belongs to the chapter's own grammar.

Nor can that first movement cancel what follows. At verse 51 the language
sharpens: the bread that Christ will give is his flesh for the life of the
world. The dispute turns on giving flesh to eat. Jesus does not retreat to the
earlier wording. He joins eating his flesh and drinking his blood to eternal
life, resurrection on the last day, true food and true drink, and mutual
abiding (6:53--58). The discourse holds together descent from the Father,
Incarnation, faith, sacrificial self-gift, communion, and bodily resurrection.
An interpretation that preserves only one of these movements has not preserved
the chapter.

Verse 63 in modern numbering---verse 64 in the
Douay--Rheims/Vulgate---``It is the spirit that quickeneth: the flesh profiteth
nothing''---does not identify Christ's flesh as useless. Within the same
discourse, that would make verses 51--58 self-refuting. In context, the
contrast is best read between divine gift received through the Spirit and a
merely natural or carnal apprehension that measures Christ by ordinary
physical possibilities. John Chrysostom makes the
same distinction when he insists that ``flesh'' in verse 63 names carnal
hearing, not the flesh Christ has just said is necessary for life.

The withdrawal of many disciples matters. John does not stage a misunderstanding
followed by an explanatory correction accepted by all. He records division,
departure, Peter's confession, and Judas's shadow (6:60--71). The reader is not
licensed to make the saying easier than the Evangelist makes it. Neither is the
reader licensed to imagine crude dismemberment: the giver is the incarnate Son
who will pass through death and raise the eater on the last day.

